\chapter{Considerações Finais}
\label{sec:conclusao}

Este capítulo apresenta as considerações finais do trabalho, sintetizando os principais pontos discutidos ao longo do documento, bem como as contribuições esperadas, limitações identificadas e o planejamento das atividades futuras da pesquisa.

\section{Síntese do Trabalho}
\label{subsec:sintese}

\textbf{Ideias a abordar:}
\begin{itemize}
    \item Relembrar o problema:
    \begin{itemize}
        \item uso de LLMs para preparação PMP
        \item incerteza quanto à confiabilidade
    \end{itemize}
    \item Reforçar o objetivo do estudo
    \item Retomar a abordagem proposta:
    \begin{itemize}
        \item avaliação empírica baseada em questões PMP
        \item uso de diferentes estratégias de prompting
    \end{itemize}
    \item Destacar os resultados preliminares:
    \begin{itemize}
        \item evidência de viabilidade
        \item tendências iniciais observadas
    \end{itemize}
\end{itemize}

\section{Atividades Futuras}
\label{subsec:atividades_futuras}

\textbf{Ideias a abordar:}
\begin{itemize}
    \item Expansão do dataset:
    \begin{itemize}
        \item aumento do número de questões
        \item maior cobertura dos domínios PMP
    \end{itemize}
    \item Inclusão de novos modelos de LLM
    \item Refinamento das estratégias de prompting:
    \begin{itemize}
        \item prompts mais estruturados
        \item uso de grounding em fontes do PMI
    \end{itemize}
    \item Execução de experimentos em larga escala
    \item Análise estatística:
    \begin{itemize}
        \item testes de significância
        \item comparação robusta entre cenários
    \end{itemize}
    \item Aprofundamento da análise qualitativa:
    \begin{itemize}
        \item categorização mais refinada dos erros
    \end{itemize}
    \item Escrita e submissão de artigos científicos
\end{itemize}

\section{Cronograma de Execução}
\label{subsec:cronograma}

\textbf{Ideias a abordar:}
\begin{itemize}
    \item Organização temporal das atividades restantes
    \item Distribuição por meses ou trimestres
\end{itemize}

\begin{table}[h]
\centering
\caption{Cronograma de execução da pesquisa}
\begin{tabular}{lcccc}
\hline
\textbf{Atividade} & \textbf{M1--M2} & \textbf{M3--M4} & \textbf{M5--M6} & \textbf{M7--M8} \\
\hline
Expansão do dataset & X & X &  &  \\
Refinamento de prompts & X & X &  &  \\
Execução dos experimentos &  & X & X &  \\
Análise dos resultados &  &  & X & X \\
Redação da dissertação &  &  & X & X \\
Submissão de artigos &  &  &  & X \\
\hline
\end{tabular}
\end{table}

